\documentclass[11pt,a4paper]{article}
\usepackage[T1]{fontenc}
\usepackage{lmodern,microtype}
\usepackage[margin=21mm,headheight=15pt]{geometry}
\usepackage{amsmath,amssymb,mathtools,mathpartir}
\usepackage{booktabs,tabularx,array,enumitem,xcolor}
\usepackage{fancyhdr,hyperref}
\definecolor{ink}{HTML}{173A52}
\definecolor{muted}{HTML}{52616B}
\hypersetup{colorlinks=true,linkcolor=ink,urlcolor=ink,
  pdftitle={Probabilistic Relational Outcome Logic (pRelOL): Implemented Rules},
  pdfauthor={pRelOL study checkpoint}}
\pagestyle{fancy}
\fancyhf{}
\fancyhead[L]{\small\color{muted}pRelOL}
\fancyhead[R]{\small\color{muted}Study checkpoint / prototype-3}
\fancyfoot[L]{\small\color{muted}22 September 2026}
\fancyfoot[R]{\small\thepage}
\renewcommand{\headrulewidth}{0.3pt}
\setlength{\parindent}{0pt}
\setlength{\parskip}{6pt}
\setlength{\emergencystretch}{2em}
\setlist{nosep,leftmargin=1.5em,topsep=4pt}
\renewcommand{\arraystretch}{1.25}
\newcommand{\scope}[1]{\textcolor{muted}{\small\textbf{Scope.} #1}\par\medskip}
\newcommand{\nextpage}[2]{\clearpage\section{#1}\scope{#2}}
\newcommand{\src}[1]{\par\smallskip{\small\textcolor{muted}{\textbf{Lean source:}} \path{#1}}\par}
\newcommand{\Qd}[4]{\{#1\}\;#2\sim#3\;\{#4\}}
\newcommand{\Prh}[4]{[\,#1\,]\;#2\sim#3\;[\,#4\,]}
\newcommand{\Cpl}[3]{\mathsf{Cpl}_{#1}(#2,#3)}
\newcommand{\cp}[1]{\mathsf{couple}(#1)}
\newcommand{\bind}{\mathbin{\gg\!=}}
\newcommand{\pure}{\eta}
\newcommand{\emp}{\mathsf{emp}}
\newcommand{\ex}[2]{\mathsf{exactly}(#1,#2)}
\newcommand{\mass}{\mathsf{mass}}
\newcommand{\supp}{\mathsf{supp}}
\newcommand{\while}[2]{\mathsf{while}\ #1\ \mathsf{do}\ #2}
\newcommand{\sk}{\mathsf{skip}}
\newcommand{\cond}[3]{\mathsf{if}\ #1\ \mathsf{then}\ #2\ \mathsf{else}\ #3}
\newcommand{\sample}[1]{\mathsf{sample}(#1)}
\newcommand{\assign}[1]{\mathsf{assign}(#1)}
\newcommand{\Loop}{\mathsf{Loop}}
\newcommand{\Rge}{\mathbb{R}_{\geq 0}^{\infty}}
\newcommand{\Ninf}{\mathbb{N}^{\infty}}
\newcommand{\mEq}[2]{\mathsf{massEq}(#1,#2)}
\newcommand{\rulelabel}[1]{\textsc{#1}}

\begin{document}
\begin{center}
{\LARGE\bfseries\color{ink}Probabilistic Relational}\\[3pt]
{\LARGE\bfseries\color{ink}Outcome Logic (pRelOL)}\\[5pt]
{\Large Implemented rules}\\[7pt]
{\small A reading and pencil-proof companion to the Lean development}\\[3pt]
{\small Mathematical baseline: \texttt{prototype-3}, commit \texttt{3542464}}\\[3pt]
{\small Source paths refreshed for the reorganized repository}
\end{center}

\section{The objects being related}
\scope{This handbook selects the core implemented rules. It is not an exhaustive
API listing or a claim of completeness. The companion theorem index links the
statements to their Lean proofs.}

A state relation $R(a,b)$ talks about two individual states. An \emph{outcome
assertion} $P(x,y)$ talks about two whole collections of outcomes. Keep these
levels separate throughout a proof.

\begin{center}
\begin{tabularx}{\linewidth}{@{}lX@{}}
\toprule
Symbol & Meaning \\
\midrule
$M\alpha$ & Outcomes over states of type $\alpha$. Sets and weighted functions are instances. \\
$0,\ x+y,\ \pure a$ & Empty outcome, additive combination, and pure outcome. \\
$x\bind k$ & Run a state kernel $k:\alpha\to M\gamma$ from every outcome in $x$. \\
$P\Rightarrow Q$ & $\forall x,y,\ P(x,y)\Rightarrow Q(x,y)$. \\
$P^{\mathsf T}(y,x)$ & $P(x,y)$, with the two arguments exchanged. \\
\bottomrule
\end{tabularx}
\end{center}

The common kernel judgment is
\[
 \boxed{\Qd{P}{k}{l}{Q}\ \Longleftrightarrow\
 \forall x,y,\ P(x,y)\Rightarrow Q(x\bind k,y\bind l).}
\]
For arbitrary outcome transformers $F,G$, replace the final pair by $(F(x),G(y))$;
Lean calls that judgment \path{quadSem}. For commands, use their semantic kernels.

Assertions used repeatedly below are
\begin{align*}
 \emp(x,y)&\iff x=0\ \land\ y=0,\\
 \ex{x_0}{y_0}(x,y)&\iff x=x_0\ \land\ y=y_0,\\
 (P\oplus Q)(x,y)&\iff \exists x_1,x_2,y_1,y_2,\quad
 x=x_1+x_2\ \land\ y=y_1+y_2\\[-2pt]
 &\hspace{32mm}{}\land P(x_1,y_1)\land Q(x_2,y_2).
\end{align*}
Thus $\oplus$ decomposes both outcomes. It is different from logical disjunction,
which simply asks whether $P(x,y)$ or $Q(x,y)$ holds. Pieces may have overlapping
support; their weights still add.

\textbf{How to use this packet.} Read a rule, write out its quantifiers, try its
proof on paper, and only then inspect Lean. Use \path{README.md} for the study
schedule, \path{exercises.md} for exercises, and \path{theorem-index.md} for exact
declaration names. Paths here are relative to \path{docs/study/}; source paths
below are relative to \path{RelationalOutcomeLogic/}. Module paths have changed;
Lean declaration namespaces are preserved.
\src{Generic/Assertions.lean}
\src{Generic/Judgment.lean}

\nextpage{Shared structural rules}{Any \texttt{OutcomeModel}: an additive commutative
monoid with the pure/bind monad laws, zero annihilation, and distributivity of bind
over addition in both arguments. No idempotence or normalization is assumed.}

Write $(k;k')(a)=k(a)\bind k'$ and $(k+k')(a)=k(a)+k'(a)$.
\begin{mathpar}
 \inferrule*[right=Skip]{ }{\Qd{P}{\pure}{\pure}{P}}
 \and
 \inferrule*[right=Consequence]{P'\Rightarrow P\quad \Qd{P}{k}{l}{Q}\quad Q\Rightarrow Q'}
 {\Qd{P'}{k}{l}{Q'}}
 \\
 \inferrule*[right=Sequence]{\Qd{P}{k}{l}{Q}\\\Qd{Q}{k'}{l'}{S}}
 {\Qd{P}{k;k'}{l;l'}{S}}
 \and
 \inferrule*[right=Swap]{\Qd{P}{k}{l}{Q}}
 {\Qd{P^{\mathsf T}}{l}{k}{Q^{\mathsf T}}}
 \\
 \inferrule*[right=Split]{\Qd{P_1}{k}{l}{Q_1}\\\Qd{P_2}{k}{l}{Q_2}}
 {\Qd{P_1\oplus P_2}{k}{l}{Q_1\oplus Q_2}}
 \\
 \inferrule*[right=Additive-choice]{\Qd{P}{k_1}{l_1}{Q_1}\\\Qd{P}{k_2}{l_2}{Q_2}}
 {\Qd{P}{k_1+k_2}{l_1+l_2}{Q_1\oplus Q_2}}
 \\
 \inferrule*[right=Disjunction]{\Qd{P_1}{k}{l}{Q}\\\Qd{P_2}{k}{l}{Q}}
 {\Qd{P_1\lor P_2}{k}{l}{Q}}
\end{mathpar}

\textbf{Empty outcomes and exact inputs.}
\[
 \Qd{\emp}{k}{l}{\emp},\qquad
 \Qd{P}{(\lambda a.\,0)}{(\lambda b.\,0)}{\emp}.
\]
\[
 \Qd{\ex{x}{y}}{k}{l}{Q}\iff Q(x\bind k,y\bind l).
\]
The empty assertion is satisfiable at $(0,0)$; it is different from the always-false
assertion. An empty-output kernel has no terminating outcomes.

\textbf{Why splitting works.} If $x=x_1+x_2$, bind yields
$(x_1\bind k)+(x_2\bind k)$. These two actual output pieces witness the conclusion.
For arbitrary transformers $F,G$, this rule needs their additivity as an explicit
premise; it does not follow merely from their type.

\textbf{Why choice has a different precondition.} Both branch pairs start from
the same $x,y$. Their outputs are then added. Over probability weights, adding
two normalized kernels has mass two: this is a rule for raw weighted kernels,
not a probabilistic command that secretly chooses a branch with probability $1/2$.

\src{Generic/Rules.lean}
\textbf{Main names:} \path{quad_seq}, \path{quad_split}, \path{quad_choice}.

\nextpage{Weights and exact couplings}{A semiring $K$ equipped with
\texttt{SummationSemiring}. Multiplication need not commute. Positivity and
absence of zero divisors are not premises of coupling composition.}

An outcome is a function $\mu:\alpha\to K$. The chosen indexed sum defines
\[
 \pure a(x)=\begin{cases}1&x=a,\\0&x\ne a,\end{cases}
 \qquad (\mu\bind k)(z)=\sum_a\mu(a)\,k(a)(z),
 \qquad \mass(\mu)=\sum_a\mu(a).
\]
Support means nonzero weight: $\supp(\mu)=\{a\mid\mu(a)\ne0\}$.

A joint $J:\alpha\times\beta\to K$ is an exact $R$-coupling of $\mu,\nu$ when
\[
 \forall a,\ \sum_b J(a,b)=\mu(a),\qquad
 \forall b,\ \sum_a J(a,b)=\nu(b),\qquad
 J(a,b)\ne0\Rightarrow R(a,b).
\]
Write $\Cpl{R}{\mu}{\nu}$ for existence of such a joint and
$\cp{R}(\mu,\nu)$ for the corresponding outcome assertion.

\begin{mathpar}
 \inferrule*[right=Pure]{R(a,b)}{\Cpl{R}{\pure a}{\pure b}}
 \and
 \inferrule*[right=Add]{\Cpl{R}{\mu_1}{\nu_1}\\\Cpl{R}{\mu_2}{\nu_2}}
 {\Cpl{R}{\mu_1+\mu_2}{\nu_1+\nu_2}}
 \\
 \inferrule*[right=Scale]{\Cpl{R}{\mu}{\nu}}
 {\Cpl{R}{r\mu}{r\nu}}
 \and
 \inferrule*[right=Weaken]{\Cpl{R}{\mu}{\nu}\\\forall a,b,\ R(a,b)\Rightarrow S(a,b)}
 {\Cpl{S}{\mu}{\nu}}
 \\
 \inferrule*[right=Bind]{\Cpl{R}{\mu}{\nu}\\
 \forall a,b,\ R(a,b)\Rightarrow\Cpl{S}{k(a)}{l(b)}}
 {\Cpl{S}{\mu\bind k}{\nu\bind l}}
\end{mathpar}
Here $r\mu$ multiplies each weight on the \emph{left}. A bind witness is built
by choosing a continuation joint $J_{a,b}$ and summing
$J(a,b)\,J_{a,b}(c,d)$ over $(a,b)$, in that multiplication order.

Two consequences are especially useful:
\[
 \Cpl{R}{\mu}{\nu}\Rightarrow\mass(\mu)=\mass(\nu),\qquad
 \Cpl{=}{\mu}{\nu}\iff\mu=\nu.
\]
For probabilities, exact couplings match termination mass as well as outcomes.
A bad cross pair does not by itself refute existence of a suitable joint.

\src{Algebra/Weighted/Coupling.lean}
\textbf{Main names:} \path{Coupled.bind}, \path{Coupled.mass_eq},
\path{coupled_eq_iff}.

\nextpage{The pRHL fragment and command rules}{The bridge is generic for weighted
kernels. The command rules here use \texttt{Prob.Cmd} with ENNReal weights.
Every sampling constructor carries a proof of mass one at every input state.}

Define the statewise judgment by
\[
 \Prh{R}{k}{l}{S}\iff
 \forall a,b,\ R(a,b)\Rightarrow\Cpl{S}{k(a)}{l(b)}.
\]
The proved semantic equivalence is
\[
 \boxed{\Prh{R}{k}{l}{S}\iff\Qd{\cp{R}}{k}{l}{\cp{S}}.}
\]
Forward: compose an input coupling with the statewise joints. Backward: apply the
outcome judgment to related pure inputs. This explains how the two interfaces
share proofs. It is not a translation of an external procedure or module language.

\begin{mathpar}
 \inferrule*[right=Assign]{\forall a,b,\ R(a,b)\Rightarrow S(f(a),g(b))}
 {\Prh{R}{\assign{f}}{\assign{g}}{S}}
 \\
 \inferrule*[right=Sample]{\forall a,\mass(k(a))=1\quad\forall b,\mass(l(b))=1\\
 \forall a,b,\ R(a,b)\Rightarrow\Cpl{S}{k(a)}{l(b)}}
 {\Prh{R}{\sample{k}}{\sample{l}}{S}}
 \\
 \inferrule*[right=If]{R\Rightarrow(g\leftrightarrow h)\\
 \Prh{R\land g}{c_1}{d_1}{S}\\\Prh{R\land\neg g}{c_2}{d_2}{S}}
 {\Prh{R}{\cond{g}{c_1}{c_2}}{\cond{h}{d_1}{d_2}}{S}}
\end{mathpar}
Here $g$ tests the left state and $h$ the right. Skip, consequence, sequence,
symmetry, and input case analysis are also implemented.

\textbf{One-sided steps.} Assignment on the left against skip requires
$R(a,b)\Rightarrow S(f(a),b)$. Left sampling against skip requires normalization
and
\[
 R(a,b)\Rightarrow\forall t\in\supp(k(a)),\ S(t,b).
\]
A left conditional against any fixed right command $d$ is proved with two
premises, under $R\land g$ and $R\land\neg g$; it needs no right guard.
Right-sided counterparts follow by symmetry.

\textbf{Quantitative assertion rule.} If $\Qd{P}{c}{d}{Q}$ then
$\Qd{r\cdot P}{c}{d}{r\cdot Q}$, where
$(r\cdot P)(\mu,\nu)\iff\exists\mu_0,\nu_0,\ \mu=r\mu_0\land\nu=r\nu_0\land P(\mu_0,\nu_0)$.
This is an image assertion, not division by $r$.

\src{Generic/Coupling.lean}
\src{Subsumption/PRHL/Embedding.lean}
\src{Subsumption/PRHL/OneSided.lean}
\src{Prob/Logic/Rules.lean}

\nextpage{Loop semantics and synchronous reasoning}{Generic iteration uses
\texttt{IterationModel} and complete lattices of outcomes. Weighted loop couplings
use \texttt{OmegaSummationSemiring}. One-sided losslessness below uses probability.}

The finite approximants count loop unfoldings, starting with no terminating output:
\begin{align*}
 L_0(s)&=0,\\
 L_{n+1}(s)&=\begin{cases}k(s)\bind L_n&g(s),\\\pure s&\neg g(s),\end{cases}\\
 \Loop(g,k)(s)&=\sup_n L_n(s).
\end{align*}
The sequence is increasing. Bind preserves its increasing limit, yielding both
the unfolding equation and the least-fixed-point characterization:
\[
 \Loop(g,k)(s)=
 \begin{cases}k(s)\bind\Loop(g,k)&g(s),\\\pure s&\neg g(s).\end{cases}
\]
An always-true loop has empty terminating semantics. Probability commands have
termination mass at most one; missing mass represents nontermination.

\begin{mathpar}
 \inferrule*[right=While-sync]{R\Rightarrow(g\leftrightarrow h)\\
 \Prh{R\land g}{k}{l}{R}}
 {\Prh{R}{\Loop(g,k)}{\Loop(h,l)}{R\land\neg g\land\neg h}}
\end{mathpar}
There is no separate termination premise. The proof builds one product execution
whose increasing joint approximants have the required marginals. Arbitrarily chosen
finite couplings need not form an increasing sequence, so they cannot simply be
passed to the limit. The outcome-level version replaces brackets by coupling
assertions using the pRHL bridge.

\textbf{A loop against skip needs more.} Let
\[
 \mathsf{LosslessOn}(A,c)\iff
 \forall a,\ A(a)\Rightarrow\mass(\mathsf{sem}(c)(a))=1.
\]
\begin{mathpar}
 \inferrule*[right=While-left]{\Prh{R\land g}{c}{\sk}{R}\\
 \mathsf{LosslessOn}(\lambda a.\exists b,\ R(a,b),\while{g}{c})}
 {\Prh{R}{\while{g}{c}}{\sk}{R\land\neg g}}
\end{mathpar}
This premise concerns the \emph{whole loop}. A lossless body is insufficient:
$\while{\mathsf{true}}{\sk}$ has mass zero, whereas skip has mass one. Exact
marginals cannot relate them. The right-sided theorem is symmetric.

\src{Generic/Loops/Semantics.lean}
\src{Generic/Loops/Synchronous.lean}
\src{Subsumption/PRHL/OneSided.lean}

\nextpage{Ranked asynchronous loops}{Ordered summation semirings; probability and
Set interfaces also exist. The schedule is supplied by the proof author. No
automatic invariant or scheduler discovery is implemented.}

Two loops need not advance together. A scheduler chooses a mode for each state
pair $p=(a,b)$ and a rank $\rho(p)\in\mathbb N$. All obligations below are required
whenever the invariant $R(a,b)$ holds and that mode is selected.

\begin{center}
\begin{tabularx}{\linewidth}{@{}l l X@{}}
\toprule
Mode & Guard obligation & Required local coupling \\
\midrule
Left & $g(a)$ & Couple $k(a)$ with $\pure b$, supported by $R$ and strict rank decrease. \\
Right & $h(b)$ & Couple $\pure a$ with $l(b)$, supported by $R$ and strict rank decrease. \\
Both & $g(a)\land h(b)$ & Couple $k(a)$ with $l(b)$, supported by $R$. \\
Done & $\neg g(a)\land\neg h(b)$ & Exit; no body step. \\
\bottomrule
\end{tabularx}
\end{center}

More precisely, for a unilateral step from $p$, the coupling relation is
\[
 R_{<p}(x,y)\iff R(x,y)\land\rho(x,y)<\rho(p).
\]
The left obligation is $\Cpl{R_{<p}}{k(a)}{\pure b}$; the right is
$\Cpl{R_{<p}}{\pure a}{l(b)}$. The both obligation is
$\Cpl{R}{k(a)}{l(b)}$.

\begin{mathpar}
 \inferrule*[right=While-async]{\text{A schedule satisfying the four obligations above}}
 {\Prh{R}{\Loop(g,k)}{\Loop(h,l)}{R\land\neg g\land\neg h}}
\end{mathpar}

\textbf{What the theorem proves.} The chosen local joints define a product loop.
Its full output joint has the \emph{original two loops} as exact marginals. These
marginal equations are a theorem, not extra promises hidden in the schedule.

\textbf{What the rank ensures.} Every supported unilateral successor lowers the
natural rank. Consequently an uninterrupted unilateral phase cannot continue
forever. Both-mode steps may reset the rank, and infinitely many joint steps may
still diverge. There is no uniform global bound on the number of loop steps.

Coupling a unilateral body with a pure state also forces that body's total weight
to be one. In probability this is losslessness of the selected step.

\textbf{Boundary to remember.} Almost-sure termination of an unbounded random
retry phase does not supply this strict support-wise rank. The existing mixed rule
does not cover such a phase merely from its probability-one termination proof.
The separate lossless loop-versus-skip rule on the previous page can handle cases
outside this scheduler's premises.

\src{Generic/Loops/Asynchronous.lean}
\textbf{Main names:} \path{AsyncSchedule}, \path{AsyncSchedule.isCoupling_loop},
\path{prhl_while_async}, \path{quad_couple_while_async}.

\nextpage{Quantitative refutation certificates}{Probability weights, or the generic
ordered summation-semiring instance. Only the selected events must be finite;
whole output distributions need not have finite support.}

For finite $A\subseteq\alpha$ and $B\subseteq\beta$, write
$\mu(A)=\sum_{a\in A}\mu(a)$. Define the cover condition by
\[
 \mathsf{Cover}(R,A,B)\iff
 \forall a\in A,\ \forall b,\ R(a,b)\Rightarrow b\in B.
\]
The quantifier over $b$ is global, including states outside the support of $\nu$.
Every exact coupling obeys the event bound
\[
 \Cpl{R}{\mu}{\nu}\ \land\ \mathsf{Cover}(R,A,B)
 \quad\Rightarrow\quad \mu(A)\leq\nu(B).
\]
Hence a \texttt{MassDeficitCert} stores $A,B$, their cover proof, and a strict gap:
\begin{mathpar}
 \inferrule*[right=Deficit]{\mathsf{Cover}(R,A,B)\\\nu(B)<\mu(A)}
 {\neg\Cpl{R}{\mu}{\nu}}
 \and
 \inferrule*[right=Bounds]{\mathsf{Cover}(R,A,B)\\
 L\leq\mu(A)\quad\nu(B)\leq U\quad U<L}
 {\neg\Cpl{R}{\mu}{\nu}}
\end{mathpar}
Two lower bounds cannot establish a deficit. In particular, finite loop unrolling
usually gives a lower bound on terminating outcomes, not the right-side upper
bound required here. No subtraction or division is used in the certificate.

\textbf{Refuting the actual contract.} Let $F,G$ be the two actual outcome
transformers. A violating pair of permitted inputs refutes the judgment:
\begin{mathpar}
 \inferrule*[right=Refute]{P(x,y)\\Q\Rightarrow\cp{R}\\
 \neg\Cpl{R}{F(x)}{G(y)}}
 {\neg\Qd{P}{F}{G}{Q}}
\end{mathpar}
Here the displayed judgment is the transformer version, \texttt{quadSem}.
The program certificate \texttt{BugCert} fixes states $a_0,b_0$ and a deficit on
$F(\pure a_0),G(\pure b_0)$; its use still requires the original precondition at
those pure inputs. A failure to discover a joint is not a refutation theorem.

\textbf{Concrete calculation.} Under equality, take $A=B=\{\mathsf{false}\}$.
If one output assigns mass $1/2$ to false and the other assigns $1/4$, the cover
is immediate and the strict gap refutes equality coupling. Matching supports
do not repair the mismatch. This is the shape of the cached-randomness bug.

\src{Prob/Refutation/Certificate.lean}
\src{Algebra/Weighted/Events.lean}
\src{Generic/Refutation.lean}
\textbf{Main names:} \path{MassDeficitCert.refutes},
\path{MassDeficitCert.of_bounds}, \path{BugCert.refutes_of_entails}.

\nextpage{Observations, splitting, and the search contract}{The rich observation
and mass-labelled splitting APIs here are implemented in the probability layer.
Restrictions keep the original weights; they do not renormalize.}

An observation $f:\alpha\to\gamma$ maps outputs by
$f_*\mu=\mu\bind(\lambda a.\pure(f(a)))$. If
$R(a,b)\Rightarrow S(f(a),g(b))$, then
\[
 \Cpl{R}{\mu}{\nu}\Rightarrow\Cpl{S}{f_*\mu}{g_*\nu}.
\]
A deficit on these \emph{actual observed outputs} refutes an internal contract $Q$
when $Q(\mu,\nu)\Rightarrow\Cpl{S}{f_*\mu}{g_*\nu}$. One must also prove the
certificate input satisfies the original precondition.

Write $\mEq{p}{q}(\mu,\nu)$ for $\mass(\mu)=p\land\mass(\nu)=q$.
\textbf{Split one joint by a shared event.} Suppose $\Cpl{R}{\mu}{\nu}$ and
$R(a,b)\Rightarrow(E(a)\leftrightarrow D(b))$. Write $\mu|_E$ for restriction,
$p=\mass(\mu|_E)$, and $q=\mass(\mu|_{\neg E})$. Then
\[
 \Bigl((\cp{R\land E\land D}\land\mEq{p}{p})\ \oplus\
 (\cp{R\land\neg E\land\neg D}\land\mEq{q}{q})\Bigr)(\mu,\nu).
\]
The two pieces come from restricting the \emph{same joint}. Their marginals are
exactly the corresponding restrictions, even when one piece has mass zero.

\textbf{The implemented end-to-end example.} A trial has miss weight $q$ and hit
weight $r$, with $q+r=1$. Both inputs are pure $(n,\mathsf{false})$: budget $n$
and no hit yet. Full scan performs all trials; early exit stops at the first hit. Define
\[
 S_0=0,\qquad S_{n+1}=qS_n+r,\qquad q^n+S_n=1.
\]
The contract is a split into the following two outcome classes:
\begin{center}
\begin{tabularx}{\linewidth}{@{}l l X@{}}
\toprule
Piece & Mass on each side & Coupled relation \\
\midrule
Failure & $q^n$ & Both answers are false; both remaining budgets are zero; each used $n$ trials. \\
Success & $S_n$ & Both answers are true; early exit uses at most as many trials as full scan; full scan used $n$. \\
\bottomrule
\end{tabularx}
\end{center}
\path{search_outcome_quad} proves this contract using asynchronous alignment and
joint splitting. For $n\geq2$ and $r>0$, strict trial savings have positive
probability. These are consumed-trial results, not an expected-runtime theorem.

The faulty cached implementation samples once and reuses the answer. For
$0<q<1$ and $n>1$, its failure probability is $q>q^n$. Observe only the answer
flag, place cached output on the deficit's left, and transfer the contradiction
back to the original contract. \path{cached_outcome_spec_refuted} refutes precisely
the contract proved for early exit.

\src{Prob/Coupling/Observation.lean}
\src{Prob/Refutation/Observation.lean}
\src{Prob/Coupling/Restriction.lean}
\src{Examples/Probability/Search/Case.lean}

\nextpage{The Set layer and the family loop rule}{Set outcomes use union as
addition. This page records Set-specific APIs; the general weighted
outcome-assertion family loop rule has not been implemented.}

For sets $X,Y$, distinguish the following assertions:
\begin{align*}
 \mathsf{atomW}(R)(X,Y)&\iff\forall a\in X,\ \forall b\in Y,\ R(a,b),\\
 \mathsf{atom}(R)(X,Y)&\iff X\ne\varnothing\land Y\ne\varnothing
                  \land\mathsf{atomW}(R)(X,Y),\\
 \mathsf{coupleL}(R)(X,Y)&\iff\forall a\in X,\ \exists b\in Y,\ R(a,b),\\
 \cp{R}(X,Y)&\iff\mathsf{coupleL}(R)(X,Y)
          \land\forall b\in Y,\ \exists a\in X,\ R(a,b).
\end{align*}
Bitotal coupling has a witness relation whose two projections are exactly $X,Y$.
It can hold for two empty sets. Adding nonemptiness gives
$\mathsf{coupleNE}(R)$, so
\[
 (\mathsf{coupleNE}(R)\oplus\top)(X,Y)
 \iff\exists a\in X,\ \exists b\in Y,\ R(a,b).
\]
Without nonemptiness, $\cp{R}\oplus\top$ is always true by selecting an empty
coupling piece. Refuting a product atom may use a bad reachable cross pair;
refuting a coupling needs to rule out possible partners, not just one pairing.

\textbf{Family loop rule.} Let $\phi_n$ describe active outcomes after $n$ body
steps and $\psi_n$ outcomes that exit then. Require:
\begin{enumerate}
 \item $\phi_n(X,Y)$ forces every state of $X$ to satisfy $g$ and every state of $Y$ to satisfy $h$.
 \item $\psi_n(X,Y)$ forces every state of $X$ to satisfy $\neg g$ and every state of $Y$ to satisfy $\neg h$.
 \item For every $n$, $\Qd{\phi_n}{c}{d}{\phi_{n+1}\oplus\psi_{n+1}}$.
 \item For \emph{every} family $E_n,E'_n$, if $\psi_n(E_n,E'_n)$ for all $n$, then
 $\psi_\infty(\bigcup_n E_n,\bigcup_n E'_n)$. Lean calls this \texttt{Converges}.
\end{enumerate}
Then
\[
 \boxed{\Qd{\phi_0\oplus\psi_0}{\while{g}{c}}{\while{h}{d}}{\psi_\infty}.}
\]
There is no termination premise. All terminating exits are collected by the union.
The convergence premise quantifies over arbitrary families satisfying $\psi_n$;
it does not secretly assume these families are increasing.

\textbf{One-sided choice.} From $\Qd{P}{c_1}{d}{Q_1}$ and
$\Qd{P}{c_2}{d}{Q_2}$, Set RelOL derives
$\Qd{P}{c_1+c_2}{d}{Q_1\oplus Q_2}$. This duplicates the passive outcome using
$Y\cup Y=Y$. The same argument fails for non-idempotent probability addition.

\src{Set/Assertions.lean}
\src{Set/Coupling.lean}
\src{Set/Rules.lean}

\nextpage{Shared assumptions and the five scalar models}{The scalar construction
uses total summation semirings. A bare semiring provides neither arbitrary indexed
sums nor the order and continuity needed for full loops.}

The assumptions have separate jobs:
\begin{center}
\begin{tabularx}{\linewidth}{@{}lX@{}}
\toprule
Interface & What it provides \\
\midrule
\texttt{OutcomeModel} & Empty, addition, pure, bind; monoid, monad, and distributive laws. Shared structural rules. \\
\texttt{SummationSemiring} & Chosen arbitrary indexed sums with singleton, addition, exchange, product-index, and left/right multiplication laws. Constructs weighted bind and couplings. \\
\texttt{PositiveSummation} & A sum is zero iff all summands are zero. With no zero divisors, supplies exact support descriptions. \\
\texttt{IterationModel} & Complete lattices of outcomes, empty as bottom, monotone bind, and preservation of increasing countable limits in the continuation. \\
\texttt{OmegaSummationSemiring} & Complete scalar lattice with zero as bottom, monotone operations, and increasing-limit laws for sums and multiplication. Derives weighted iteration. \\
\bottomrule
\end{tabularx}
\end{center}

All five implemented scalar models satisfy the required summation and increasing
limit laws, as well as positivity and absence of zero divisors.
\begin{center}
\begin{tabularx}{\linewidth}{@{}l l l l X@{}}
\toprule
Model & $+$ & $\cdot$ & $0,1$ & Meaning \\
\midrule
Reachability & $\lor$ & $\land$ & false, true & Possible output \\
$\Rge$ & $+$ & $\times$ & $0,1$ & Nonnegative weight \\
$\Ninf$ & $+$ & $\times$ & $0,1$ & Trace count \\
Tropical cost & $\min$ & $+$ & $\infty,0$ & Cheapest cost \\
Word languages & $\cup$ & concat. & $\varnothing,\{\epsilon\}$ & Emitted finite words \\
\bottomrule
\end{tabularx}
\end{center}
For tropical costs the semantic order reverses numerical order, so semantic
suprema become numerical infima. Language multiplication is noncommutative;
the order of emissions matters. Languages may contain infinitely many finite words.
ENNReal weights become probabilities only when
the appropriate mass constraints are established.

\textbf{One calculation, several meanings.} A direct route of weight $d$ or two
successive edges of weights $a,b$ gives output weight $d+ab$. Boolean weights
record reachability; unit counting weights give two traces; costs $7,2,3$ give
minimum cost $5$. A weighted countdown gives $r^n\pure 0$ in all instances.

\textbf{Scope.} This is not every partial algebra from Outcome Logic: operations
are total, weights have no countable-support field, iteration uses increasing
sequential limits, and trace words are finite. Compatibility theorems connect
the generic construction to the existing Set and ENNReal developments.

{\small\textcolor{muted}{\textbf{Lean sources:}}
\path{Algebra/Scalar/Basic.lean}, \path{Algebra/Scalar/Order.lean},
\path{Bridges/Set.lean}, \path{Bridges/Probability.lean},
\path{Examples/Algebra/Routes.lean}.}

\nextpage{A checkpoint before further development}{Use this page to test
understanding after the exercises. The aim is to reconstruct why a proof is valid,
not to memorize theorem names or read every file in order.}

\textbf{Reconstruct these arguments without opening their proofs.}
\begin{enumerate}[itemsep=7pt]
 \item \textbf{Sequence and split.} Explain exactly where associativity and
 additivity of bind enter. Give the witness decomposition for split.
 \item \textbf{The pRHL bridge.} Prove both directions. Identify where related
 pure inputs are used and where arbitrary coupled inputs are used.
 \item \textbf{Coupling versus product atoms.} Give a finite example where an
 equality coupling exists although a cross pair fails equality.
 \item \textbf{Deficit certificate.} Calculate the two-coin event masses,
 prove global cover, and state the original precondition and postcondition
 implication needed to refute a program judgment.
 \item \textbf{Synchronous limits.} Explain why a coherent increasing sequence
 of joints is sufficient and unrelated finite witnesses are insufficient.
 \item \textbf{Asynchronous alignment.} Draw a schedule for a machine with a
 bookkeeping phase. Show the guard and rank obligation in each mode.
 \item \textbf{Search.} Derive the budget-two result by hand, then explain the
 arbitrary-budget proof using alignment, restriction, and a mass recurrence.
 \item \textbf{Generalization.} Reinterpret one weighted calculation as
 reachability, counting, cost, and a trace language. Identify the scalar laws used.
\end{enumerate}

\bigskip
\textbf{What this checkpoint already contains.}
\begin{itemize}[itemsep=3pt]
 \item Set and weighted outcome assertions, exact couplings, and refutation.
 \item Semantic exact-coupling pRHL equivalence and core command rules.
 \item Full loop semantics, synchronous rules, lossless one-sided rules, and
 ranked mixed asynchronous alignment.
 \item The full-scan/early-exit contract and a cached-randomness refutation of
 that same contract; shared algebra with five scalar instances and examples.
\end{itemize}

\bigskip
\textbf{What remains a future direction.} Informed symbolic execution and
coupling/certificate discovery, richer case studies, a general weighted
outcome-family loop rule, broader asynchronous progress criteria, and external
LeanCrypt integration. The current Set \texttt{symexec} is a simplification wrapper.
These are not needed to complete the study checkpoint. Completeness is not a goal.

\textbf{Reading question.} When comparing Outcome Logic, pRHL/product programs,
and Quantitative Weakest Hyper Pre, distinguish semantic expressiveness from
local proof rules, reusable invariants, automation, and empirical proof effort.
This development does not establish a new expressiveness separation from QWHP.

\textbf{Ready to resume?} Explain the search correctness and bug proofs from the
definitions, reconstruct several small proofs independently, and list the exact
assumptions of the next rule or tactic you want. Use the companion study plan
for selected paper readings and a place to record unresolved questions.

\end{document}
